\documentclass[10pt,a4paper]{amsart}
\usepackage{amsmath,amssymb,amsthm,graphicx}
\usepackage{verbatim}
\newenvironment{reference}{\comment}{\endcomment}
%\newenvironment{reference}{}{}
\newenvironment{slogan}{\comment}{\endcomment}
\newenvironment{history}{\comment}{\endcomment}

% For commutative diagrams we use Xy-pic
\usepackage[all]{xy}

% We use 2cell for 2-commutative diagrams.
\xyoption{2cell}
% Theorem environments.
%
\theoremstyle{plain}
\newtheorem{theorem}[subsection]{Theorem}
\newtheorem{proposition}[subsection]{Proposition}
\newtheorem{lemma}[subsection]{Lemma}

\theoremstyle{definition}
\newtheorem{definition}[subsection]{Definition}
\newtheorem{example}[subsection]{Example}
\newtheorem{exercise}[subsection]{Exercise}
\newtheorem{situation}[subsection]{Situation}

\theoremstyle{remark}
\newtheorem{remark}[subsection]{Remark}
\newtheorem{remarks}[subsection]{Remarks}

\numberwithin{equation}{subsection}

% Macros
%
\def\lim{\mathop{\mathrm{lim}}\nolimits}
\def\colim{\mathop{\mathrm{colim}}\nolimits}
\def\Spec{\mathop{\mathrm{Spec}}}
\def\Hom{\mathop{\mathrm{Hom}}\nolimits}
\def\Ext{\mathop{\mathrm{Ext}}\nolimits}
\def\SheafHom{\mathop{\mathcal{H}\!\mathit{om}}\nolimits}
\def\SheafExt{\mathop{\mathcal{E}\!\mathit{xt}}\nolimits}
\def\Sch{\mathit{Sch}}
\def\Mor{\mathop{\mathrm{Mor}}\nolimits}
\def\Ob{\mathop{\mathrm{Ob}}\nolimits}
\def\Sh{\mathop{\mathit{Sh}}\nolimits}
\def\NL{\mathop{N\!L}\nolimits}
\def\CH{\mathop{\mathrm{CH}}\nolimits}
\def\proetale{{pro\text{-}\acute{e}tale}}
\def\etale{{\acute{e}tale}}
\def\QCoh{\mathit{QCoh}}
\def\Ker{\mathop{\mathrm{Ker}}}
\def\Im{\mathop{\mathrm{Im}}}
\def\Coker{\mathop{\mathrm{Coker}}}
\def\Coim{\mathop{\mathrm{Coim}}}

% Boxtimes
%
\DeclareMathSymbol{\boxtimes}{\mathbin}{AMSa}{"02}

%
% Macros for moduli stacks/spaces
%
\def\QCohstack{\mathcal{QC}\!\mathit{oh}}
\def\Cohstack{\mathcal{C}\!\mathit{oh}}
\def\Spacesstack{\mathcal{S}\!\mathit{paces}}
\def\Quotfunctor{\mathrm{Quot}}
\def\Hilbfunctor{\mathrm{Hilb}}
\def\Curvesstack{\mathcal{C}\!\mathit{urves}}
\def\Polarizedstack{\mathcal{P}\!\mathit{olarized}}
\def\Complexesstack{\mathcal{C}\!\mathit{omplexes}}
% \Pic is the operator that assigns to X its picard group, usage \Pic(X)
% \Picardstack_{X/B} denotes the Picard stack of X over B
% \Picardfunctor_{X/B} denotes the Picard functor of X over B
\def\Pic{\mathop{\mathrm{Pic}}\nolimits}
\def\Picardstack{\mathcal{P}\!\mathit{ic}}
\def\Picardfunctor{\mathrm{Pic}}
\def\Deformationcategory{\mathcal{D}\!\mathit{ef}}
\usepackage{cleveref}
\crefname{originalequation}{original equation}{original equations}
\Crefname{originalequation}{Original equation}{Original equations}
\makeatletter
\newlabel{equation-list}{{048S}{source}{}{sourcepointer-1}{}}
\newlabel{equation-list@cref}{{[originalequation][0][]048S}{[1][0][]source}{}{}{}}
\makeatother
\crefname{originallemma}{original lemma}{original lemmas}
\Crefname{originallemma}{Original lemma}{Original lemmas}
\makeatletter
\newlabel{spaces-morphisms-lemma-base-change-finite-type}{{03XH}{source}{}{sourcepointer-2}{}}
\newlabel{spaces-morphisms-lemma-base-change-finite-type@cref}{{[originallemma][0][]03XH}{[1][0][]source}{}{}{}}
\makeatother
\crefname{originallemma}{original lemma}{original lemmas}
\Crefname{originallemma}{Original lemma}{Original lemmas}
\makeatletter
\newlabel{spaces-morphisms-lemma-large-enough}{{0488}{source}{}{sourcepointer-3}{}}
\newlabel{spaces-morphisms-lemma-large-enough@cref}{{[originallemma][0][]0488}{[1][0][]source}{}{}{}}
\makeatother
\crefname{originallemma}{original lemma}{original lemmas}
\Crefname{originallemma}{Original lemma}{Original lemmas}
\makeatletter
\newlabel{spaces-morphisms-lemma-locally-finite-type-surjective-geometric-points}{{0487}{source}{}{sourcepointer-4}{}}
\newlabel{spaces-morphisms-lemma-locally-finite-type-surjective-geometric-points@cref}{{[originallemma][0][]0487}{[1][0][]source}{}{}{}}
\makeatother
\crefname{originallemma}{original lemma}{original lemmas}
\Crefname{originallemma}{Original lemma}{Original lemmas}
\makeatletter
\newlabel{spaces-morphisms-lemma-permanence-finite-type}{{0462}{source}{}{sourcepointer-5}{}}
\newlabel{spaces-morphisms-lemma-permanence-finite-type@cref}{{[originallemma][0][]0462}{[1][0][]source}{}{}{}}
\makeatother
\title{Geometric point criterion for set-theoretic pre-equivalence relations}
\author{The Stacks Project Authors}
\date{stacks-project:groupoids-quotients.tex:local-snapshot-20261012; GNU FDL 1.2 or later}
\begin{document}
\maketitle
\noindent\textit{Editorial note.} All delivered mathematical text is the author's own, copied from the pinned Stacks Project chapter file groupoids-quotients.tex; no statement or proof step was rewritten, reordered or removed.\par\medskip
\noindent\textit{Editorial note.} The counted proof is the author's complete argument: The two details-omitted spots each reuse an argument displayed in full earlier in the same proof; no new technique is hidden.\par\medskip
\noindent\textit{Editorial note.} The label lemma-set-theoretic-pre-equivalence-geometric names the counted statement, and the proof environment immediately following it in the source is delivered verbatim as its proof.\par\medskip
\noindent\textit{Editorial note.} The counted statement is carried by the author's own lemma environment in the chapter source.\par\medskip
\noindent\textit{Allowed standard prerequisites (editorial).}\begin{itemize}
\item Chapter Definition definition-set-theoretic-equivalence, included as a context span below: j is a set-theoretic pre-equivalence relation when for every algebraically closed field k over B the relation sim-R on U(k), given by u sim-R u-prime iff some field extension K over k admits a K-valued point r of R with s(r) = u and t(r) = u-prime, is an equivalence relation; j is a set-theoretic equivalence relation when it is moreover universally injective
\item The three induced maps of Equation equation-list: the composition map j-comp from R composed with itself over U, the diagonal map j-diag from U, and the flip map j-flip from R, all landing in U times over B U
\item Morphisms of Spaces, Lemma spaces-morphisms-lemma-large-enough: surjectivity statements on topological spaces can be witnessed on k-valued points as soon as the cardinality of the algebraically closed field k exceeds the bound lambda associated to R
\item Morphisms of Spaces, Lemmas spaces-morphisms-lemma-permanence-finite-type and spaces-morphisms-lemma-base-change-finite-type: local finite type is preserved under composition and under base change
\item Morphisms of Spaces, Lemma spaces-morphisms-lemma-locally-finite-type-surjective-geometric-points: a nonempty algebraic space locally of finite type over an algebraically closed field has a k-valued point
\end{itemize}\medskip
\ifdefined\hypertarget\hypertarget{sourcepointer-1}{}\fi
\noindent\textit{Source reference.} Original equation 048S: sources/stacks-groupoids-quotients/pinned-original-source.tex, source lines 813--813. Prerequisite pointer invoked at 1097; it names a standard fact of the same chapter that the counted argument uses, and every step of the counted proof built on it is retained verbatim in the delivered body.\par\medskip
\ifdefined\hypertarget\hypertarget{sourcepointer-2}{}\fi
\noindent\textit{Source reference.} Original lemma 03XH: sources/stacks-chow/tags.txt, source lines 5148--5148. Prerequisite pointer invoked at 1214; it names a standard fact of a sibling chapter that the counted argument uses, and every step of the counted proof built on it is retained verbatim in the delivered body.\par\medskip
\ifdefined\hypertarget\hypertarget{sourcepointer-3}{}\fi
\noindent\textit{Source reference.} Original lemma 0488: sources/stacks-chow/tags.txt, source lines 5535--5535. Prerequisite pointer invoked at 1179; it names a standard fact of a sibling chapter that the counted argument uses, and every step of the counted proof built on it is retained verbatim in the delivered body.\par\medskip
\ifdefined\hypertarget\hypertarget{sourcepointer-4}{}\fi
\noindent\textit{Source reference.} Original lemma 0487: sources/stacks-chow/tags.txt, source lines 5534--5534. Prerequisite pointer invoked at 1217; it names a standard fact of a sibling chapter that the counted argument uses, and every step of the counted proof built on it is retained verbatim in the delivered body.\par\medskip
\ifdefined\hypertarget\hypertarget{sourcepointer-5}{}\fi
\noindent\textit{Source reference.} Original lemma 0462: sources/stacks-chow/tags.txt, source lines 5457--5457. Prerequisite pointer invoked at 1213; it names a standard fact of a sibling chapter that the counted argument uses, and every step of the counted proof built on it is retained verbatim in the delivered body.\par\medskip
\begin{definition}
\label{definition-set-theoretic-equivalence}
Let $S$ be a scheme, and let $B$ be an algebraic space over $S$.
Let $j : R \to U \times_B U$ be a pre-relation over $B$.
\begin{enumerate}
\item We say $j$ is a {\it set-theoretic pre-equivalence relation} if
for all algebraically closed fields $k$ over $B$ the relation
$\sim_R$ on $U(k)$ defined by
$$
\overline{u} \sim_R \overline{u}'
\Leftrightarrow
\begin{matrix}
\exists\text{ field extension }K/k, \ \exists\ r \in R(K), \\
s(r) = \overline{u}, \ t(r) = \overline{u}'
\end{matrix}
$$
is an equivalence relation.
\item We say $j$ is a {\it set-theoretic equivalence relation}
if $j$ is universally injective and a set-theoretic pre-equivalence
relation.
\end{enumerate}
\end{definition}


\begin{lemma}
\label{lemma-set-theoretic-pre-equivalence-geometric}
In the situation of Definition \ref{definition-set-theoretic-equivalence}.
The following are equivalent:
\begin{enumerate}
\item The morphism $j$ is a set-theoretic pre-equivalence relation.
\item The subset $j(|R|) \subset |U \times_B U|$ contains the image of
$|j'|$ for any of the morphisms $j'$ as in Equation (\ref{equation-list}).
\item For every algebraically closed field $k$ over $B$ of sufficiently large
cardinality the subset $j(R(k)) \subset U(k) \times U(k)$ is an equivalence
relation.
\end{enumerate}
If $s, t$ are locally of finite type these are also equivalent to
\begin{enumerate}
\item[(4)] For every algebraically closed field $k$ over $B$
the subset $j(R(k)) \subset U(k) \times U(k)$ is an equivalence relation.
\end{enumerate}
\end{lemma}

\begin{proof}
Assume (2). Let $k$ be an algebraically closed field over $B$.
We are going to show that $\sim_R$ is an equivalence relation.
Suppose that $\overline{u}_i : \Spec(k) \to U$, $i = 1, 2$
are $k$-valued points of $U$. Suppose that $(\overline{u}_1, \overline{u}_2)$
is the image of a $K$-valued point $r \in R(K)$. Consider the
solid commutative diagram
$$
\xymatrix{
\Spec(K') \ar@{..>}[r] \ar@{..>}[d]
&
\Spec(k) \ar[d]_{(\overline{u}_2, \overline{u}_1)} &
\Spec(K) \ar[d] \ar[l] \\
R \ar[r]^-j &
U \times_B U &
R \ar[l]_-{j_{flip}}
}
$$
We also denote $r \in |R|$ the image of $r$.
By assumption the image of $|j_{flip}|$ is contained in the image of
$|j|$, in other words there exists a $r' \in |R|$ such that
$|j|(r') = |j_{flip}|(r)$. But note that $(\overline{u}_2, \overline{u}_1)$
is in the equivalence class that defines $|j|(r')$ (by the commutativity
of the solid part of the diagram). This means there exists a field
extension $K'/k$ and a morphism $r' : \Spec(K) \to R$
(abusively denoted $r'$ as well) with
$j \circ r' = (\overline{u}_2, \overline{u}_1) \circ i$
where $i : \Spec(K') \to \Spec(K)$ is the obvious map.
In other words the dotted part of the diagram commutes.
This proves that $\sim_R$ is a symmetric relation on $U(k)$.
In the similar way, using that the image of $|j_{diag}|$ is contained
in the image of $|j|$ we see that $\sim_R$ is reflexive (details omitted).

\medskip\noindent
To show that $\sim_R$ is transitive assume given
$\overline{u}_i : \Spec(k) \to U$, $i = 1, 2, 3$
and field extensions $K_i/k$ and points
$r_i : \Spec(K_i) \to R$, $i = 1, 2$ such that
$j(r_1) = (\overline{u}_1, \overline{u}_2)$ and
$j(r_1) = (\overline{u}_2, \overline{u}_3)$. Then we may choose a
commutative diagram of fields
$$
\xymatrix{
K & K_2 \ar[l] \\
K_1 \ar[u] & k \ar[l] \ar[u]
}
$$
and we may think of $r_1, r_2 \in R(K)$. We consider the
commutative solid diagram
$$
\xymatrix{
\Spec(K') \ar@{..>}[r] \ar@{..>}[d]
&
\Spec(k) \ar[d]_{(\overline{u}_1, \overline{u}_3)} &
\Spec(K) \ar[d]^{(r_1, r_2)} \ar[l]
\\
R \ar[r]^-j &
U \times_B U &
R \times_{s, U, t} R \ar[l]_-{j_{comp}}
}
$$
By exactly the same reasoning as in the first part of the proof, but
this time using that $|j_{comp}|((r_1, r_2))$ is in the image of $|j|$,
we conclude that a field $K'$ and dotted arrows exist making the
diagram commute. This proves that $\sim_R$ is transitive and concludes
the proof that (2) implies (1).

\medskip\noindent
Assume (1) and let $k$ be an algebraically closed field over $B$ whose
cardinality is larger than $\lambda(R)$, see
Morphisms of Spaces, Lemma \ref{spaces-morphisms-lemma-large-enough}.
Suppose that $\overline{u} \sim_R \overline{u}'$ with
$\overline{u}, \overline{u}' \in U(k)$. By assumption there exists
a point in $|R|$ mapping to $(\overline{u}, \overline{u}') \in |U \times_B U|$.
Hence by
Morphisms of Spaces, Lemma \ref{spaces-morphisms-lemma-large-enough}
we conclude there exists an $\overline{r} \in R(k)$ with
$j(\overline{r}) = (\overline{u}, \overline{u}')$. In this way we see
that (1) implies (3).

\medskip\noindent
Assume (3). Let us show that
$\Im(|j_{comp}|) \subset \Im(|j|)$. Pick any point
$c \in |R \times_{s, U, t} R|$. We may represent this by a morphism
$\overline{c} : \Spec(k) \to R \times_{s, U, t} R$, with $k$ over $B$
having sufficiently large cardinality. By assumption we see that
$j_{comp}(\overline{c}) \in U(k) \times U(k) = (U \times_B U)(k)$
is also the image $j(\overline{r})$ for some $\overline{r} \in R(k)$.
Hence $j_{comp}(c) = j(r)$ in $|U \times_B U|$ as desired (with
$r \in |R|$ the equivalence class of $\overline{r}$). The same argument
shows also that $\Im(|j_{diag}|) \subset \Im(|j|)$ and
$\Im(|j_{flip}|) \subset \Im(|j|)$ (details omitted).
In this way we see that (3) implies (2). At this point we have
shown that (1), (2) and (3) are all equivalent.

\medskip\noindent
It is clear that (4) implies (3) (without any assumptions on $s$, $t$).
To finish the proof of the lemma we show that (1) implies (4) if $s, t$
are locally of finite type. Namely, let $k$ be an algebraically closed
field over $B$. Suppose that $\overline{u} \sim_R \overline{u}'$ with
$\overline{u}, \overline{u}' \in U(k)$. By assumption the algebraic space
$Z = R \times_{j, U \times_B U, (\overline{u}, \overline{u}')} \Spec(k)$
is nonempty. On the other hand, since $j = (t, s)$ is locally of finite type
the morphism $Z \to \Spec(k)$ is locally of finite type as well (use
Morphisms of Spaces, Lemmas \ref{spaces-morphisms-lemma-permanence-finite-type}
and \ref{spaces-morphisms-lemma-base-change-finite-type}).
Hence $Z$ has a $k$ point by
Morphisms of Spaces, Lemma
\ref{spaces-morphisms-lemma-locally-finite-type-surjective-geometric-points}
and we conclude that $(\overline{u}, \overline{u}') \in j(R(k))$
as desired. This finishes the proof of the lemma.
\end{proof}

\end{document}
